\fontsize{12pt}{13}\selectfont %Ocupa letra 14 de TImes New Roman

La predicción temprana de fallas de software permite optimizar el uso de recursos y tiempo durante las etapas de prueba y mantenimiento. Los modelos de aprendizaje automático supervisado empleados para esta tarea se entrenan con conjuntos de datos de métricas de código que presentan dos problemas persistentes: alta dimensionalidad y desequilibrio de clases. Esta tesis evalúa si un preprocesamiento en dos etapas ---selección de características mediante análisis de relevancia y control de redundancia, seguida de balanceo de clases--- mejora la predicción de fallas sobre BugHunter, un conjunto de datos de 15 proyectos Java con métricas en tres niveles de granularidad (archivo, clase y método).

Se aplican Ganancia de Información, Chi-Cuadrado y ReliefF para el análisis de relevancia; Incertidumbre Simétrica y Semejanza de Coseno para el control de redundancia; y Random Over-Sampling, Random Under-Sampling y SMOTE para el balanceo del conjunto de entrenamiento. Los modelos se evalúan con Random Forest en 13.916 ejecuciones, complementadas por un análisis de sensibilidad de 28.440 ejecuciones con Random Forest, XGBoost y SVM lineal que registra métricas por clase.

La selección de características reduce la mediana del número de métricas en un 40\,\%, 88\,\% y 50\,\% en los niveles de método, clase y archivo, respectivamente, con una pérdida de \textit{F-measure} inferior a 0,02, y las métricas de tamaño de código (CLOC, LOC y LLOC) son las seleccionadas con mayor frecuencia. El hallazgo más relevante concierne al balanceo: aplicado solo al entrenamiento, no mejora el \textit{F-measure} ponderado cuando el conjunto de prueba conserva la distribución real de clases (tamaño de efecto de Cliff $|\delta| < 0{,}15$), pero sí mejora de forma consistente la detección de la clase defectuosa (\textit{F-measure} de la clase \textit{bug} y MCC) en los tres clasificadores, un beneficio que la métrica agregada enmascara. Entrenar por proyecto individual supera al entrenamiento con el conjunto unificado de proyectos. La tesis revisa además el estado del arte de los modelos de lenguaje de gran escala (LLM) en la predicción de defectos y evalúa su aplicación mediante un piloto exploratorio: tres LLM aplicados directamente sobre las métricas no superan a Random Forest en MCC y reproducen la misma dependencia de la métrica. Los resultados indican que la conclusión sobre la utilidad del balanceo depende críticamente de la métrica de evaluación, una advertencia metodológica que se extiende a los enfoques basados en modelos de lenguaje.

\hfill \break
\textbf{Palabras clave:} predicción de fallas de software; selección de características; desequilibrio de clases; métricas de código; BugHunter; coeficiente de correlación de Matthews; modelos de lenguaje de gran escala
